\documentclass[11pt,a4paper]{article}
\usepackage[margin=1in]{geometry}
\usepackage{hyperref}
\usepackage{enumitem}
\usepackage{graphicx}
\usepackage{xcolor}

% -----------------------------------------------------
% bvraghav-tiet-ucs503p-article.sty
% -----------------------------------------------------

% Variable: Lab Instructor
\makeatletter \newcommand{\BvrSetLabInstructor}[1]{%
  \gdef\Bvr@LabInstructor{#1}}
% default
\newcommand{\Bvr@LabInstructor}{Dr. Jeelani %
  Asif}
% public accessor
\newcommand{\BvrLabInstructorValue}{\Bvr@LabInstructor}
\makeatother

% Add subtitle
\usepackage{titling}
% -- Set title bold
\pretitle{\begin{center}\bfseries\fontsize{18bp}{18bp}\selectfont}
\makeatletter
\newcommand{\BvrSubtitle}[1]{%
  \posttitle{%
    \par\end{center}
    \begin{center}\large#1\end{center}
    \vskip 0.5em}%
}
\makeatother

% Hints in the section title(s)
\newcommand{\BvrHint}[1]{%
  {\normalfont\normalsize\itshape\hfill #1}}

% -----------------------------------------------------

\title{Conflict Resolution System for a University Campus}

\BvrSetLabInstructor{Dr. Jeelani Asif}

\BvrSubtitle{%
  Project Proposal for UCS503P  \\
  Submitted to: \BvrLabInstructorValue \\ \bfseries
  Thapar Institute of Engineering and Technology% 
}

\author{
  Arhana Mor, CSED%
  \thanks{Roll No: \texttt{1024030773}, %
    Email: \texttt{amor\_be24\@thapar.edu}}
  \and 
  Garv Bansal, CSED%
  \thanks{Roll No: \texttt{1024030988}, %
    Email: \texttt{gbansal\_be24\@thapar.edu}}
}

\date{\today}

\begin{document}
\maketitle

\section{Higher-order goal%
  \BvrHint{\ldots secondary framing}}
This project aims to streamline university workflows and improve learning by reducing the time and effort required to allocate campus resources, including classrooms, Digi Spaces, and laboratories. The immediate engineering objective is to develop a measurable and deployable software solution for 
 \emph{automated conflict detection and resolution in venue bookings.} 

\section{Time-to-value %
\BvrHint{\ldots primary heuristic}}
\begin{itemize}[leftmargin=0.7cm]
    \item We will accelerate value delivery by rapidly prototyping the core workflow and validating it with end users within each short development cycle.
    \item We will favor continuous feedback over premature optimization, using weekly iterations and early CI-backed automation to minimize deployment overhead.
    \item Initial success will be measured through tangible, testable outcomes from the first iteration, with subsequent releases driven by validated user feedback.
\end{itemize}



\section{Problem Statement}
Campus lecture halls and laboratories are often reallocated for placement activities, orientations, and other institutional events without timely notification to the affected batches and faculty. This results in last-minute class cancellations and inefficient venue reassignment due to the lack of readily available alternatives. Common operational pain points include:
\begin{itemize}[leftmargin=0.7cm]
    \item \textbf{Delayed resource conflict notification:} Faculty and students are often informed late, or not informed at all, when their allocated classrooms or laboratories are preempted.
    \item \textbf{No alternative venue recommendations:} The current process does not always identify and suggest suitable replacement rooms based on capacity, availability, or facility requirements.
    \item \textbf{Manual, communication-dependent workflow:} Venue changes rely heavily on emails and manual coordination, increasing the risk of missed updates, inconsistent information, and delayed resolution.
    \item \textbf{Limited allocation transparency:} Students and faculty lack a centralized view of venue availability, booking changes, and the status of resource conflicts.
\end{itemize}

\textbf{Why this matters now (contextual relevance):} During placement examinations for pre-final and final-year students, campus resources can remain occupied until 8 PM, creating conflicts with scheduled laboratory sessions. Many of these disruptions could be mitigated through intelligent reassignment—for example, relocating a lab session to an available classroom in Blocks D, E, or F while allowing students to use their own laptops.

\section{Proposed Solution %
\BvrHint{\ldots Software Proposition}}
\subsection{Overview}
Build a \textbf{web-based} “Conflict-to-Resolution” system with the following components:
\begin{itemize}[leftmargin=0.7cm]
    \item \textbf{Venue booking portal:} Timetable coordinators submit booking requests specifying required capacity, facilities (projectors, computers, etc.), time slot, and activity type (lecture, placement examination, orientation, etc.).
    \item \textbf{Automated conflict detection:} The system validates requests against existing bookings and generates alerts when resource conflicts or double bookings are detected.
    \item \textbf{Alternative venue recommendation:} For conflicting requests, the system identifies suitable available venues based on capacity, required facilities, and time constraints.
    \item \textbf{Conflict resolution workflow:} Coordinators can review the conflicting activities and select an appropriate action, such as relocating or rescheduling a lower-priority activity.
    \item \textbf{Automated notifications:} Affected students and faculty are notified immediately when a scheduled activity is relocated, suspended, or otherwise modified.
\end{itemize}

\subsection{Core workflow}
\begin{enumerate}[leftmargin=0.7cm]
    \item The coordinator defines the venue requirements, including seating capacity, required resources, time slot, and activity type, or directly specifies a preferred room (e.g., LT402).
    \item The system evaluates available rooms against the specified constraints and ranks suitable venues based on resource and scheduling compatibility.
    \item The coordinator selects a preferred venue from the available options.
    \item If the selected venue is already occupied, the system displays the conflicting activity, expected occupancy, required resources, and responsible faculty member or supervisor.
    \item The coordinator determines which activity retains the venue based on factors such as activity priority, timing, and resource requirements.
    \item For the displaced activity, the system identifies and recommends currently unoccupied venues that satisfy its capacity, resource, and scheduling requirements.
    \item The coordinator selects an alternative venue or cancels/reschedules the affected activity. The system records the resolution and notifies the relevant students and faculty.
\end{enumerate}



\subsection{Operational constraints for a University Campus}
\begin{itemize}[leftmargin=0.7cm]
    \item Ensure the web application remains responsive and functional on low-to-mid-range devices and variable campus network conditions.
    \item Prioritize deterministic conflict detection, constraint-based venue matching, and workflow automation over computationally intensive or research-oriented algorithms.
    \item Use widely supported web technologies, standard hosting infrastructure, and managed databases to enable cost-effective deployment and maintenance.
    \item Design the system to handle concurrent booking requests and updates while maintaining consistent room availability and allocation state.
\end{itemize}

\section{Solution Approach %
\BvrHint{\ldots Engineering focus}}
This engineering project prioritizes scalable administrative workflows and reliable system automation over the development of research-oriented techniques.

\subsection{Resource Recommendations %
\BvrHint{\ldots non-research}}
Recommendations will use a rule-based ranking approach:
\begin{itemize}[leftmargin=0.7cm]
    \item Match required constraints such as capacity, facilities, activity type, and time slot against eligible resources.
    \item Rank unoccupied rooms above occupied rooms. For occupied rooms, display basic conflict details such as the current activity, batch, faculty/supervisor, and occupancy duration.
    \item For a conflict, recommend alternative rooms for both activities. If no suitable unoccupied room exists, consider relaxed alternatives based on factors such as greater distance or reduced resource availability while preserving essential constraints.
    \item Present the top-ranked alternatives to the coordinator for final selection; the system will not automatically relocate or cancel activities.
\end{itemize}

\subsection{Web architecture %
\BvrHint{\ldots web-based solution}}
A typical 3-tier web architecture:
\begin{itemize}[leftmargin=0.7cm]
    \item \textbf{Frontend:} Responsive interface for coordinators, faculty, and students, including venue booking, conflict views, recommendations, and notifications.
    \item \textbf{Backend API:} Handles authentication, booking requests, conflict detection, room recommendation, allocation updates, and notification workflows.
    \item \textbf{Database:} Stores users, rooms, resource specifications, bookings, conflict records, activity details, and allocation history.
\end{itemize}

\subsection{CI/CD and fast delivery enabler}
CI/CD will be used to:
\begin{itemize}[leftmargin=0.7cm]
    \item Automatically build and execute tests for every code change.
    \item Generate consistent, reproducible deployment artifacts for staging.
    \item Support frequent and reliable releases of incremental features.
\end{itemize}

\section{Evaluation Criterion %
\BvrHint{\ldots measurable and attributable}}
We will evaluate using quantitative metrics directly tied to the core venue-allocation workflow.

\subsection{Primary evaluation metric}
\textbf{Alternative Recommendation Validity:} Percentage of system-generated recommendations that satisfy all mandatory constraints and are available for the requested time slot.
\begin{itemize}[leftmargin=0.7cm]
    \item Target: $\geq$ 95\% (calculated using (valid recommendations / total recommendations) × 100)
    \item Attribution: Validity is verified against the booking and room-resource data stored in the database.
\end{itemize}

\subsection{Secondary evaluation metrics}
\begin{itemize}[leftmargin=0.7cm]
    \item \textbf{Recommendation Proximity:} distance between the original venue and the recommended alternative, measured using predefined campus coordinates.
    \item \textbf{Conflict Resolution Coverage:} Percentage of detected conflicts for which the system generates at least one feasible alternative for either affected activity.
    \item \textbf{Notification Delivery Rate:} Percentage of affected users for whom a notification delivery event is successfully recorded.
    \item \textbf{Workflow Efficiency:} Number of manual actions required by the coordinator to identify a feasible alternative and record the resolution.
\end{itemize}

\subsection{Pilot validation plan %
\BvrHint{\ldots high-level}}
\begin{itemize}[leftmargin=0.7cm]
    \item Conduct a pilot with 1 coordinator, 1 lecturer affected by a venue conflict, and 10–15 student participants to validate the deployed prototype.
    \item Evaluate the core workflow, including conflict detection, alternative-room recommendations, booking updates, and notification delivery.
    \item Check whether affected users receive and correctly view notifications following a venue change or cancellation.
\end{itemize}


\section{Scalability %
\BvrHint{\ldots with foresight}}
The proposed system is designed to scale both in terms of system load and practical deployment requirements.

\begin{enumerate}[leftmargin=0.7cm]
    \item \textbf{Scalable (theoretically):} The platform should support increasing numbers of venues, bookings, users, and concurrent requests through:
    \begin{itemize}[leftmargin=0.7cm]
        \item decoupled frontend and backend services,
        \item database indexing and pagination for frequent booking and availability queries,
        \item using stateless API design.
    \end{itemize}
    
    \item \textbf{Operationally deployable (practically):} The system should be deployable on standard cloud or web-hosting infrastructure using:
    \begin{itemize}[leftmargin=0.7cm]
        \item managed relational database services
        \item containerized or platform-hosted deployment,
        \item CI/CD pipelines for automated staging and production releases.
    \end{itemize}
\end{enumerate}

\section{Engine availability heuristic}
A mature software stack is available to implement the proposed campus resource allocation platform:
\begin{itemize}[leftmargin=0.7cm]
    \item Web framework: REST APIs and responsive frontend interfaces.
    \item Database: Managed relational database for users, venues, bookings, and allocation history.
    \item Authentication: Session- or token-based authentication with role-based access control.
    \item Notification service: Notifications for booking changes and conflict resolution.
    \item Testing framework: Unit and integration testing for workflow and allocation logic.
    \item CI/CD infrastructure: Automated build, testing, and staging deployment pipeline.
\end{itemize}
These established components allow the project to focus on workflow automation, allocation correctness, and measurable performance rather than developing the infrastructure from scratch.

\section{Project scope and deliverables %
\BvrHint{\ldots iteration-friendly}}
\subsection{Initial deliverable %
\BvrHint{\ldots first iteration}}
\begin{itemize}[leftmargin=0.7cm]
    \item Venue requirement form + Conflict detection
    \item Post-conflict recommendations: System recommends only currently unoccupied rooms that satisfy the affected activity's requirements.
    \item Public announcements: Display cancellation or relocation notices for affected classes, without requiring student authentication in the initial prototype.
    \item CI/CD: Automated build, backend testing, and deployment to a staging environment.
\end{itemize}

\subsection{Subsequent deliverables}
\begin{itemize}[leftmargin=0.7cm]
    \item Incorporate venue proximity into the ranking of alternative rooms.
    \item Notification layer: Add automated in-app notifications for affected students and faculty.
    \item Introduce authenticated access.
    \item Enhanced recommendation logic: Introduce weighted ranking based on proximity, resource availability, capacity, and other operational constraints.
\end{itemize}
\section{Risks and mitigations}
\begin{itemize}[leftmargin=0.7cm]
    \item \textbf{Unauthorized access:} Implement role-based access control so only authorized coordinators can create, modify, relocate, or cancel venue bookings. Students should have read-only access to relevant announcements and schedule changes.
    \item \textbf{Excessive data collection:} Store only operationally necessary data such as booking details, venue requirements, user roles, and allocation history.
    \item \textbf{Evaluation ambiguity:} Log booking, conflict, recommendation, and resolution events with timestamps and define clear state transitions before implementation.
    \item \textbf{Pilot deployment issues:} Deploy to staging early and use CI/CD to ensure reproducible builds, automated testing, and consistent deployments.
\end{itemize}

\section{Summary}
This project proposes a deployable web platform for efficient campus resource allocation and conflict resolution. It meets the course criteria by emphasizing time-to-value through rapid iterations, adopting CI/CD for frequent and reliable releases, and evaluating performance using measurable metrics such as Conflict Resolution Coverage. The project focuses on practical workflow automation, system reliability, and scalability rather than research-driven algorithmic novelty.

\end{document}
